\vspace{-1mm}
\section{Conclusion}
\label{sec:conclusion}
\vspace{-2mm}
We presented State of Thought (SoT), a reasoning paradigm that changes the control variable of test-time reasoning from external token programs to endogenous state-conditioned evidence organization. Rather than prescribing longer chains, larger search budgets, or fixed memory rules, SoT uses the model's evolving representational state to dynamically gate which evidence remains active and when reasoning should terminate.  This closed-loop formulation yields measurable gains on the accuracy–efficiency frontier across text and multimodal benchmarks, and generalizes robustly to limited-access and trajectory-level evaluation settings. Our results point to a broader principle: reasoning systems need not rely on increasingly elaborate external scaffolding — adaptivity and efficiency can emerge directly from the model's internal state geometry.